% This chapter is the complete working manuscript for the sixth seminar talk.
% The final section contains supplementary material outside the default live
% route.

\chapter{Derived bordism and fundamental classes}
\label[chapter]{chap:Derived_Bordism_And_Fundamental_Classes}

The preceding chapter constructed a derived zero locus
\[
  Z^{\mathrm{der}}(s)
  =
  M\mathbin{\times_E^{\mathrm{der}}}M
\]
for every section $s\colon M\to E$ of a finite-rank vector bundle, without
requiring $s$ to be transverse to the zero section. It then calculated this
object through a Koszul presentation. We now ask a different question. If $M$
is compact and $s'$ is a transverse perturbation of $s$, how are the derived
zero locus $Z^{\mathrm{der}}(s)$ and the ordinary manifold $Z(s')$ related?

They are generally not equivalent. They need not even have homeomorphic
underlying spaces. The relation between them is deliberately coarser:

\begin{quote}
  A compact derived zero locus and a transverse perturbation represent the
  same bordism class.
\end{quote}

This statement gives derived intersections their first topological payoff.
The unperturbed intersection exists as a geometric object, while its bordism
class is the classical class obtained by perturbation. The same statement also
shows how much bordism forgets. A derived object may have nontrivial animated
structure even when its bordism class vanishes.

There is one scope convention which will remain in force throughout the
chapter. Spivak's bordism theorem concerns compact \emph{quasi-smooth} derived
manifolds, meaning the derived manifolds locally presented as zero loci of
maps between ordinary manifolds. All derived intersections and vector-bundle
zero loci used in this chapter are of this kind. We do not claim the theorem
for arbitrary higher-amplitude objects in the full $\infty$-category
$\mathrm{DMfd}$. Carchedi--Steffens compare Spivak's framework with the modern
one used in this manuscript, so the quasi-smooth objects considered below may
be regarded as objects of the same $\infty$-category as in the preceding
chapters \cite[Theorem~5.10]{Carchedi_Steffens_Universal_Derived_Manifolds}.

\section{From perturbations to derived cobordisms}
\label[section]{sec:From_Perturbations_To_Derived_Cobordisms}

\subsection{A homotopy of sections}
\label[subsection]{sec:Homotopy_Of_Sections_Cobordism}

Let $M$ be a compact manifold, let $q\colon E\to M$ be a rank-$r$ vector
bundle, and let $s\colon M\to E$ be a section. Recall that its derived zero
locus is defined by the pullback square
\[
\begin{tikzcd}
  Z^{\mathrm{der}}(s) \rar \dar \drar[pullback]
    & M \dar["0"] \\
  M \rar["s"'] & E .
\end{tikzcd}
\]
Its underlying set is the ordinary zero set of $s$, but its structure sheaf
retains the equations defining that set.

Choose another section $s'\colon M\to E$. We first assume that $s'$ is
transverse to the zero section. Choose a smooth function
\[
  \beta\colon\R\longrightarrow[0,1]
\]
which is identically zero on a neighborhood of $0$ and identically one on a
neighborhood of $1$. Fiberwise affine combination defines a section
\[
  H\colon\R\times M\longrightarrow\operatorname{pr}_M^*E,
  \qquad
  H(u,m)=(1-\beta(u))s(m)+\beta(u)s'(m).
\]
Its derived zero locus is taken over $\R\times M$:
\[
\begin{tikzcd}
  W \rar \dar \drar[pullback]
    & \R\times M \dar["0"] \\
  \R\times M \rar["H"'] & \operatorname{pr}_M^*E .
\end{tikzcd}
\]
Projection to the first factor induces a map
\[
  p\colon W=Z^{\mathrm{der}}(H)\longrightarrow\R.
\]
The important point is that the single derived object $W$ contains both zero
loci as fibers.

\begin{proposition}[A homotopy of sections gives a cobordism]
  \label[proposition]{prop:Homotopy_Of_Sections_Derived_Cobordism}
  With notation as above, the map $p\colon W\to\R$ is proper and is a product
  near its fibers over $0$ and $1$. Moreover,
  \[
    W_{p=0}\simeq Z^{\mathrm{der}}(s),
    \qquad
    W_{p=1}\simeq Z(s').
  \]
\end{proposition}

\begin{proof}
  Pulling the defining square of $W$ back along $\{0\}\to\R$ replaces $H$
  by $s$, so
  \[
    W_{p=0}\simeq Z^{\mathrm{der}}(s).
  \]
  The same argument at $1$ gives
  $W_{p=1}\simeq Z^{\mathrm{der}}(s')$. Since $s'$ is transverse to the zero
  section, its derived zero locus is the ordinary zero manifold $Z(s')$.

  Choose $\varepsilon>0$ so that $\beta$ is zero on
  $(-\varepsilon,\varepsilon)$ and one on
  $(1-\varepsilon,1+\varepsilon)$. Over these two intervals the section $H$
  is constant in the parameter. Consequently,
  \[
    W|_{(-\varepsilon,\varepsilon)}
    \simeq
    Z^{\mathrm{der}}(s)\times(-\varepsilon,\varepsilon)
  \]
  and
  \[
    W|_{(1-\varepsilon,1+\varepsilon)}
    \simeq
    Z(s')\times(1-\varepsilon,1+\varepsilon).
  \]

  It remains to check properness. Let $K\subset\R$ be compact. The
  underlying space of $p^{-1}(K)$ is the ordinary zero set of $H$ inside
  $K\times M$. It is closed because the zero section of the vector bundle is
  closed. Since $K\times M$ is compact, the space underlying $p^{-1}(K)$ is
  compact. Thus $p$ is proper.
\end{proof}

The proof separates the two roles of the hypotheses. The constancy of the
homotopy near $0$ and $1$ supplies product neighborhoods of the endpoint
fibers. Compactness of $M$ prevents zeros from escaping to infinity while the
section is perturbed.

\subsection{The definition of derived cobordism}
\label[subsection]{sec:Definition_Derived_Cobordism}

The preceding construction contains every clause in the general definition.

\begin{definition}[Collared map]
  \label[definition]{def:Collared_Map_Derived_Manifolds}
  Let $p\colon W\to\R$ be a map of derived manifolds and let $i\in\R$. We say
  that $p$ is \emph{$i$-collared} if there is an $\varepsilon>0$ and an
  equivalence over $(i-\varepsilon,i+\varepsilon)$
  \[
    W|_{|p-i|<\varepsilon}
    \simeq
    W_{p=i}\times(-\varepsilon,\varepsilon).
  \]
  On the right, the interval is translated so that its origin maps to $i$.
\end{definition}

\begin{definition}[Derived cobordism]
  \label[definition]{def:Derived_Cobordism}
  Let $X_0$ and $X_1$ be compact quasi-smooth derived manifolds. A
  \emph{derived cobordism} from $X_0$ to $X_1$ consists of a quasi-smooth
  derived manifold $W$ and a proper map
  \[
    p\colon W\longrightarrow\R
  \]
  such that:
  \begin{enumerate}
    \item The map $p$ is $0$-collared and $1$-collared.
    \item There are equivalences
      \[
        W_{p=0}\simeq X_0,
        \qquad
        W_{p=1}\simeq X_1.
      \]
  \end{enumerate}
  If the objects carry maps to a manifold $T$, a derived cobordism
  \emph{over $T$} also carries a continuous map from the underlying space of
  $W$ to $T$ which restricts to the endpoint maps.
\end{definition}

This is Spivak's boundaryless formulation of cobordism
\cite[Definitions~3.5--3.6]{Spivak_Derived_Smooth_Manifolds}. Instead of
introducing derived manifolds with boundary, one uses a manifold without
boundary equipped with a proper map to $\R$. Properness is the compactness
condition appropriate to such a family. The collars permit two cobordisms to
be glued along a common endpoint fiber.

\begin{remark}[Cobordism is not equivalence]
  \label[remark]{rem:Cobordism_Not_Equivalence}
  A homotopy of sections does not generally produce an equivalence between
  their derived zero loci. It produces a larger derived object over a
  parameter line. Passing from objects to bordism classes discards the
  difference between the two endpoint fibers.
\end{remark}

\begin{corollary}[Homotopic sections]
  \label[corollary]{cor:Homotopic_Sections_Derived_Cobordant}
  Let $M$ be compact and let $s_0,s_1$ be sections of the same vector bundle
  over $M$. Then their derived zero loci are derived cobordant. If $s_1$ is
  transverse, then $Z^{\mathrm{der}}(s_0)$ is derived cobordant to the
  ordinary manifold $Z(s_1)$.
\end{corollary}

\begin{proof}
  Apply the construction above to a homotopy of sections which is constant
  near its endpoints. Transversality is needed only to identify the second
  endpoint with an ordinary manifold.
\end{proof}

\section{Classical and derived bordism}
\label[section]{sec:Classical_And_Derived_Bordism}

\subsection{The comparison map}
\label[subsection]{sec:Derived_Bordism_Comparison_Map}

Let $T$ be a manifold. Write $\Omega_d(T)$ for the ordinary unoriented
bordism group of proper maps from $d$-manifolds to $T$. Equivalently,
$\Omega_*(T)$ is the homology theory represented by $MO$. Define
$\Omega_d^{\mathrm{der}}(T)$ in the same way using compact quasi-smooth
derived manifolds of virtual dimension $d$ and derived cobordisms over $T$.
Disjoint union gives the group operation.

The inclusion of ordinary manifolds into derived manifolds induces a map
\[
  i_*
  \colon
  \Omega_*(T)\longrightarrow\Omega_*^{\mathrm{der}}(T).
\]
At first sight the target might be larger. Derived manifolds provide many new
representatives, including singular zero loci and nontransverse
intersections. The comparison theorem says that these representatives do not
produce new bordism classes.

\begin{theorem}[Spivak's bordism comparison]
  \label[theorem]{thm:Spivak_Derived_Bordism_Comparison}
  For every manifold $T$, inclusion induces an isomorphism
  \[
    i_*
    \colon
    \Omega_*(T)
    \iso
    \Omega_*^{\mathrm{der}}(T).
  \]
\end{theorem}

The absolute assertion is Spivak's Theorem 3.12, and the relative assertion is
Corollary 3.13
\cite[Theorem~3.12 and Corollary~3.13]{Spivak_Derived_Smooth_Manifolds}.
The theorem says that every compact quasi-smooth derived manifold is
\emph{derived cobordant} to an ordinary manifold. It does not say that it is
equivalent to one.

\subsection{Global equations for a compact derived manifold}
\label[subsection]{sec:Global_Equations_Compact_Derived_Manifold}

The proof reduces the general case to the construction from the first
section. Two geometric inputs make this possible.

\begin{theorem}[Embedding and global zero-locus presentation]
  \label[theorem]{thm:Compact_Derived_Manifold_Global_Zero_Locus}
  Let $X$ be a compact quasi-smooth derived manifold in Spivak's framework.
  Then there are:
  \begin{enumerate}
    \item An embedding $j\colon X\to\R^N$.
    \item An open neighborhood $U\subseteq\R^N$ of the underlying space of
      $X$.
    \item A finite-rank vector bundle $E\to U$.
    \item A section $s\colon U\to E$.
  \end{enumerate}
  With this data, there is an equivalence over $U$
  \[
    X\simeq Z^{\mathrm{der}}(s).
  \]
\end{theorem}

\begin{proof}[Source and proof architecture]
  The embedding is Spivak's Proposition 3.3. Once $X$ is embedded, the
  normal-bundle axiom, Definition 2.1(7), supplies the neighborhood $U$, the
  vector bundle, and the defining section
  \cite[Definition~2.1(7) and Proposition~3.3]
  {Spivak_Derived_Smooth_Manifolds}. The proof of the embedding theorem is a
  derived version of the familiar partition-of-unity proof for compact
  manifolds. We use the result as a quoted geometric input.
\end{proof}

There is a small but essential difference from the opening construction. The
ambient manifold $U$ need not be compact. Compactness of $X$ nevertheless
allows the perturbation to be supported near its zero locus.

Choose a compact subset $K\subset U$ whose interior contains the underlying
space of $X$. Since $s$ has no zeros outside the interior of $K$, relative
transversality gives a transverse section $s'$ which agrees with $s$ outside
$K$. The straight-line homotopy, made constant near its endpoints, is also
fixed outside $K$. Every zero of every section in the homotopy therefore lies
in $K$, which restores the properness argument.

\subsection{Proof of the comparison theorem}
\label[subsection]{sec:Proof_Derived_Bordism_Comparison}

\begin{proof}[Proof of \Cref{thm:Spivak_Derived_Bordism_Comparison}]
  We first prove surjectivity in the absolute case. Let $X$ be a compact
  quasi-smooth derived manifold. Choose a global zero-locus presentation
  \[
    X\simeq Z^{\mathrm{der}}(s)
  \]
  as in \Cref{thm:Compact_Derived_Manifold_Global_Zero_Locus}. Perturb $s$ to
  a transverse section $s'$ while keeping the perturbation fixed outside a
  compact neighborhood of the zero locus. The homotopy-zero-locus
  construction of \Cref{prop:Homotopy_Of_Sections_Derived_Cobordism} gives a
  derived cobordism
  \[
    X\sim_{\mathrm{bord}}Z(s').
  \]
  The endpoint $Z(s')$ is an ordinary compact manifold. Hence every derived
  bordism class is in the image of $i_*$.

  For injectivity, suppose that ordinary compact manifolds $M_0$ and $M_1$
  are derived cobordant. Let
  \[
    p\colon W\longrightarrow\R
  \]
  be a derived cobordism between them. Restrict to the part of $W$ over an
  open interval containing $[0,1]$. Spivak's relative embedding theorem
  embeds this restricted cobordism over the parameter into
  $\R^N\times\R$ \cite[Corollary~3.4]{Spivak_Derived_Smooth_Manifolds}.
  The normal-bundle axiom presents it as the derived zero locus of a section.
  Because $p$ is collared, this presentation is product-like near $M_0$ and
  $M_1$. Relative transversality permits us to perturb the defining section
  while keeping those collared regions fixed. The perturbed zero locus is an
  ordinary smooth cobordism from $M_0$ to $M_1$. Thus two ordinary manifolds
  which become derived cobordant were already classically cobordant, and
  $i_*$ is injective.

  For the relative theorem over $T$, one performs the same constructions
  while retaining the map to $T$. The only additional point is that a map
  from the compact zero locus to the manifold $T$ extends to a neighborhood
  in the ambient Euclidean space. The perturbation and cobordism can then be
  carried out over $T$.
\end{proof}

The two halves of the proof have distinct meanings. Surjectivity says that
derived representatives can be regularized without changing their bordism
class. Injectivity says that allowing derived cobordisms does not identify two
classical bordism classes which were previously distinct.

\begin{example}[A singular fiber of a Morse function]
  \label[example]{ex:Figure_Eight_Derived_Bordism}
  Let $f\colon T^2\to\R$ be a Morse function and let $c$ be a critical value
  of index one. The underlying space of the derived fiber
  \[
    T^2\mathbin{\times_{\R}^{\mathrm{der}}}\{c\}
  \]
  is a figure eight. It is not the underlying space of a one-dimensional
  manifold. Moving $c$ slightly to either side produces a regular fiber,
  either one circle or two circles depending on the side. The parameter
  interval gives derived cobordisms from the singular fiber to both regular
  fibers. Thus the two regular fibers are classically cobordant, while the
  singular derived fiber remains a distinct geometric object. This is one of
  Spivak's examples \cite[Example~2.7]{Spivak_Derived_Smooth_Manifolds}.
\end{example}

\section{Intersection products and Euler classes}
\label[section]{sec:Intersection_Products_And_Euler_Classes}

\subsection{The nontransverse cup-product formula}
\label[subsection]{sec:Nontransverse_Cup_Product_Formula}

Let $M$ be a closed manifold, and let $A,B\subseteq M$ be closed
submanifolds. We use cohomological notation for their Poincare-dual
unoriented bordism classes:
\[
  [A]\in MO^{\operatorname{codim}(A)}(M),
  \qquad
  [B]\in MO^{\operatorname{codim}(B)}(M).
\]
Equivalently, one may work homologically with the proper maps $A\to M$ and
$B\to M$ and then apply bordism Poincare duality. This convention prevents a
common ambiguity: the intersection formula is a cup-product formula in
cohomology, while the geometric representatives are manifolds over $M$.

If $A$ and $B$ are transverse, their ordinary intersection represents the
cup product. Derived geometry removes the transversality condition from the
geometric representative.

\begin{theorem}[General cup-product formula]
  \label[theorem]{thm:General_Cup_Product_Derived_Intersection}
  Let $A$ and $B$ be closed submanifolds of a closed manifold $M$. Then
  \[
    [A]\smile[B]
    =
    [A\mathbin{\times_M^{\mathrm{der}}}B]
  \]
  in the unoriented bordism cohomology of $M$.
\end{theorem}

\begin{proof}
  Perturb the inclusion $A\to M$ through smooth maps until it is transverse
  to $B\to M$. Denote the perturbed map by $A'\to M$. Pulling the homotopy
  back along $B\to M$ gives a derived cobordism over $M$ from
  \[
    A\mathbin{\times_M^{\mathrm{der}}}B
  \]
  to the ordinary transverse fiber product $A'\pitchfork B$. Hence the two
  objects determine the same class over $M$. The classical transverse
  intersection formula gives
  \[
    [A]\smile[B]=[A'\pitchfork B].
  \]
  Combining the two equalities proves the result.
\end{proof}

This is Spivak's Corollary 3.13
\cite[Corollary~3.13]{Spivak_Derived_Smooth_Manifolds}. Notice where
transversality remains in the argument. It is not needed to form the derived
intersection. It is used only to identify the bordism class of that canonical
object with the familiar classical product.

\subsection{The self-intersection of a zero section}
\label[subsection]{sec:Zero_Section_Self_Intersection}

Let $M$ be a closed $m$-manifold and let $q\colon E\to M$ be a rank-$r$
vector bundle. Write $z\colon M\to E$ for the zero section. Its derived
self-intersection is
\[
  X_E
  =
  M\mathbin{\times_E^{\mathrm{der}}}M.
\]
The two maps to $E$ in this expression are both $z$. The underlying space of
$X_E$ is all of $M$, but the derived object remembers that the same normal
equation has been imposed twice.

Choose a transverse section $s'\colon M\to E$. Its graph is a perturbation of
one copy of the zero section. The transverse intersection of $z(M)$ and
$s'(M)$ is precisely the zero manifold $Z(s')$. The homotopy from $z$ to
$s'$ therefore gives a derived cobordism over $M$
\[
  X_E
  \sim_{\mathrm{bord}}
  Z(s').
\]

The bordism Euler class
\[
  e_{MO}(E)\in MO^r(M)
\]
is represented by the zero locus of a transverse section. We have proved the
following self-intersection formula.

\begin{corollary}[Zero-section self-intersection]
  \label[corollary]{cor:Zero_Section_Self_Intersection_Euler_Class}
  The cohomological bordism class represented by the derived
  self-intersection of the zero section is the Euler class:
  \[
    [M\mathbin{\times_E^{\mathrm{der}}}M]
    =
    e_{MO}(E).
  \]
  In homological notation, the corresponding class over $M$ is
  \[
    e_{MO}(E)\cap[M]
    =
    [Z(s')\longrightarrow M]
    \in MO_{m-r}(M).
  \]
\end{corollary}

The derived self-intersection is canonical, whereas the transverse section
$s'$ is a choice. Any two choices are joined by a homotopy of sections, so
their zero manifolds are cobordant. Derived geometry packages the class before
the choice of perturbation is made.

\begin{example}[The trivial line bundle]
  \label[example]{ex:Trivial_Line_Derived_Self_Intersection}
  Let $E=M\times\R$ be the trivial line bundle. Its derived zero-section
  self-intersection is
  \[
    X_{\underline{\R}}
    =
    M\mathbin{\times_{M\times\R}^{\mathrm{der}}}M.
  \]
  Locally, its algebra of functions is the Koszul algebra
  \[
    \bigl(\Cinfty(U)\otimes\Lambda(\eta),\partial\eta=0\bigr).
  \]
  Thus its underlying space and classical truncation are both $M$, while it
  has a nonzero degree-one class generated by $\eta$.

  On the other hand, the zero section can be perturbed to the nowhere-zero
  constant section $1$. Its zero locus is empty. Hence
  \[
    [X_{\underline{\R}}]=0
  \]
  in bordism. The derived object is visibly nonempty and nonclassical, but its
  bordism class vanishes. This is the cleanest example of the information
  which bordism deliberately forgets.
\end{example}

\section{Fundamental classes and their limits}
\label[section]{sec:Fundamental_Classes_And_Their_Limits}

\subsection{From bordism to generalized homology}
\label[subsection]{sec:From_Bordism_To_Generalized_Homology}

Let $X$ be a compact quasi-smooth derived manifold of virtual dimension $d$.
The comparison theorem supplies a class
\[
  [X]\in\Omega_d^{\mathrm{der}}
  \cong
  \Omega_d
  =MO_d.
\]
More generally, a proper map $X\to T$ supplies a class in $MO_d(T)$.

Suppose that a generalized homology theory $\mathbb E_*$ is equipped with a
morphism
\[
  MO\longrightarrow\mathbb E
\]
between the stable objects which represent the two theories. This induces a
natural transformation
\[
  MO_*(T)\longrightarrow\mathbb E_*(T).
\]

\begin{definition}[Generalized fundamental class]
  \label[definition]{def:Generalized_Fundamental_Class_Derived_Manifold}
  The image of $[X\to T]$ under this natural transformation is the
  $\mathbb E$-homology fundamental class of $X$ over $T$:
  \[
    [X\to T]_{\mathbb E}
    \in
    \mathbb E_d(T).
  \]
\end{definition}

For the live talk, the displayed morphism should be treated as the datum that
makes ordinary unoriented manifold fundamental classes available in
$\mathbb E$-homology. The point of the comparison theorem is then immediate:
the same datum assigns compatible classes to compact quasi-smooth derived
manifolds. This is the observation made in Spivak's introduction
\cite[Section~1]{Spivak_Derived_Smooth_Manifolds}.

\subsection{Orientations}
\label[subsection]{sec:Orientations_Derived_Bordism}

Everything above is unoriented. This is not a cosmetic convention. To obtain
an integral or otherwise oriented class, one must supply orientation data.
Spivak formulates an orientation using a normal bundle. A compact derived
manifold is embedded in some $\R^N$, and the normal-bundle axiom supplies a
normal vector bundle. After adding trivial summands, this normal bundle is
stable under changing the embedding. An orientation of the stable normal
bundle is the input for oriented derived bordism
\cite[Remark~3.8]{Spivak_Derived_Smooth_Manifolds}.

In the oriented case, $MO$ is replaced by the appropriate Thom object, for
example $MSO$ for oriented bordism. A morphism from that object to
$\mathbb E$ transports the oriented bordism class to $\mathbb E$-homology.
Nothing in the existence of a derived zero locus automatically provides such
an orientation.

This brief normal-bundle discussion also points forward. Bordism only retains
the stable normal information needed for a class. The next chapter will
construct a finer stable linear object intrinsic to the derived manifold: its
tangent and cotangent complexes.

\subsection{What has not been constructed}
\label[subsection]{sec:What_Derived_Bordism_Does_Not_Construct}

The phrase \enquote{fundamental class} can obscure several separate
requirements. The construction in this chapter requires a compact
quasi-smooth derived manifold, or more generally a proper map to the chosen
target. It does not prove that a moduli problem is compact. It does not add
limiting objects to a noncompact moduli space. It does not orient a determinant
line. It also does not treat isotropy or produce a virtual integration theory
for derived orbifolds or stacks.

These distinctions matter for the elliptic moduli problems later in the
seminar. A representability theorem has the form
\[
  \text{moduli functor}
  \simeq
  \text{quasi-smooth derived geometric object}.
\]
It does not by itself imply
\[
  \text{compactness},
  \qquad
  \text{orientation},
  \qquad
  \text{virtual fundamental class}.
\]
If the represented object is later shown to be compact and is equipped with
the required normal orientation, derived bordism becomes one possible route
to a class. In applications with automorphisms, compactifications, or more
general Kuranishi structures, additional theory is needed.

We can now formulate the handoff to the next chapter. The trivial-line
self-intersection in \Cref{ex:Trivial_Line_Derived_Self_Intersection} has zero
bordism class, but its derived structure has not disappeared. What intrinsic
first-order invariant retains that structure, and how does it record
deformations and obstructions? The answer will be the stable tangent complex.

\section{Summary and supplementary materials}
\label[section]{sec:Derived_Bordism_Summary_And_Supplements}

\subsection{Takeaways}
\label[subsection]{sec:Derived_Bordism_Takeaways}

The live route through the chapter can be compressed into seven statements.

\begin{enumerate}
  \item A homotopy of vector-bundle sections gives a derived family of their
    zero loci.
  \item Making the homotopy constant near its endpoints gives collars, while
    compactness or compact support gives properness.
  \item Every compact quasi-smooth derived manifold has a global zero-locus
    presentation in Spivak's framework.
  \item Inclusion induces an isomorphism from classical unoriented bordism to
    derived bordism.
  \item The theorem compares bordism classes, not derived manifolds up to
    equivalence.
  \item A nontransverse derived intersection represents the classical
    cup product, and a zero-section self-intersection represents the bordism
    Euler class.
  \item A map from $MO$ to a generalized homology theory transports the
    bordism class to a fundamental class, while compactness and orientation
    remain additional inputs.
\end{enumerate}

The remainder of the chapter is supplementary. It supplies proofs and
examples which are useful for reference but are not part of the default
70-minute route.

\subsection{Supplement: derived cobordism is an equivalence relation}
\label[subsection]{sec:Supplement_Derived_Cobordism_Equivalence_Relation}

\begin{proposition}
  \label[proposition]{prop:Derived_Cobordism_Equivalence_Relation}
  Derived cobordism is an equivalence relation on compact quasi-smooth
  derived manifolds.
\end{proposition}

\begin{proof}
  Reflexivity is represented by the projection
  \[
    X\times\R\longrightarrow\R.
  \]
  This map is proper when $X$ is compact, is collared everywhere, and has
  fiber $X$ over both $0$ and $1$.

  Symmetry follows by reversing the parameter. If $p\colon W\to\R$ is a
  cobordism from $X_0$ to $X_1$, then $1-p$ is a cobordism from $X_1$ to
  $X_0$.

  For transitivity, let $W_{01}$ be a cobordism from $X_0$ to $X_1$ and let
  $W_{12}$ be a cobordism from $X_1$ to $X_2$. Their collars identify open
  neighborhoods of the common endpoint with
  $X_1\times(-\varepsilon,\varepsilon)$. After translating the parameters,
  glue the two derived manifolds along these equivalent open subobjects. The
  product descriptions ensure that the parameter maps agree on the overlap.
  The resulting map is proper and collared at the remaining endpoints, hence
  is a derived cobordism from $X_0$ to $X_2$.
\end{proof}

This is Spivak's Proposition 3.10
\cite[Proposition~3.10]{Spivak_Derived_Smooth_Manifolds}. The collar condition
is exactly what makes the gluing step formal.

\subsection{Supplement: the relative injectivity argument}
\label[subsection]{sec:Supplement_Relative_Injectivity_Derived_Bordism}

The injectivity proof hides two relative statements. First, a derived
cobordism $p\colon W\to\R$ can be embedded over the parameter, after
restricting to a neighborhood of $[0,1]$:
\[
\begin{tikzcd}
  W' \rar[hook,"j"] \dar["p"']
    & \R^N\times(-1,2) \dar["\operatorname{pr}_2"] \\
  (-1,2) \rar[equal] & (-1,2).
\end{tikzcd}
\]
The map to the parameter remains proper. This is the role of Spivak's
Corollary 3.4.

Second, the defining section for $W'$ may be perturbed relative to the
collared regions. In those regions the endpoint fibers are already ordinary
manifolds, and the section is already transverse in the normal directions.
Relative transversality changes the section only away from those regions. Its
ordinary zero locus is therefore a smooth cobordism with the same endpoints.

For a cobordism over $T$, the map from the compact part of the zero locus to
$T$ is extended to a neighborhood before perturbing. All subsequent
constructions retain this extended map. This is why the absolute comparison
upgrades to an isomorphism of bordism homology theories, rather than merely an
isomorphism of coefficient groups.

\subsection{Supplement: the Pontryagin--Thom picture}
\label[subsection]{sec:Supplement_Pontryagin_Thom_Derived_Bordism}

For an ordinary compact manifold $N$ embedded in $\R^k$, the
Pontryagin--Thom collapse produces a stable map from a sphere to the Thom
space of the normal bundle. Stabilizing over all embeddings yields the class
of $N$ in $MO_*$. This construction explains why the normal bundle, rather
than a chosen tangent bundle, is the natural input for unoriented bordism.

For a compact quasi-smooth derived manifold $X$, the embedding and
normal-bundle theorem again supplies stable normal data. A transverse
perturbation of a defining section produces an ordinary manifold $N$ derived
cobordant to $X$. Define the Pontryagin--Thom class of $X$ to be the class of
$N$. If a second perturbation produces $N'$, a homotopy of perturbations gives
a derived cobordism from $N$ to $N'$. Injectivity in
\Cref{thm:Spivak_Derived_Bordism_Comparison} turns this into an ordinary
cobordism, so the stable class is independent of the perturbation.

This picture is useful conceptually, but it should not replace the
homotopy-of-sections construction in the live talk. The latter makes the
geometric content visible without requiring an introduction to stable
homotopy theory.

\subsection{Supplement: the Möbius line bundle}
\label[subsection]{sec:Supplement_Mobius_Euler_Class}

Let $L\to S^1$ be the Möbius line bundle. A section can be described by a
smooth function $f\colon\R\to\R$ satisfying
\[
  f(t+1)=-f(t).
\]
The function $f(t)=\sin(\pi t)$ defines a transverse section. Modulo the
identification $t\sim t+1$, it has one zero. Therefore the derived
self-intersection of the zero section is derived cobordant to one point over
$S^1$. Its mod-$2$ Euler number is nonzero.

This complements the trivial-line example. For the trivial line, a
nowhere-zero section makes the self-intersection class vanish. For the Möbius
line, every section has a zero and the transverse zero set detects the
nontrivial unoriented Euler class.

\subsection{Related literature}
\label[subsection]{sec:Derived_Bordism_Related_Literature}

The definitions of collared derived cobordism, the embedding theorem, the
comparison with ordinary bordism, and the general cup-product formula are all
from Spivak's \emph{Derived smooth manifolds}
\cite[Sections~1--3]{Spivak_Derived_Smooth_Manifolds}. The comparison which
places Spivak's model inside the universal modern framework is
Carchedi--Steffens, especially their comparison with Spivak in Section 5.2
\cite[Section~5.2]{Carchedi_Steffens_Universal_Derived_Manifolds}.

Joyce's theory of d-manifolds and d-orbifolds develops oriented bordism and
virtual classes in a framework which also accommodates finite isotropy. Those
results are relevant to later moduli problems, but they require definitions
and hypotheses not developed in this chapter. In particular, they should not
be conflated with the bordism comparison for compact quasi-smooth derived
manifolds proved here.
